\documentclass{article}
\usepackage{amsmath,amssymb,amsthm,algorithm,algorithmic}
\usepackage{hyperref}
\newtheorem{theorem}{Theorem}
\newtheorem{lemma}{Lemma}
\newtheorem{assumption}{Assumption}
\newcommand{\prox}{\operatorname{prox}}
\newcommand{\R}{\mathbb{R}}

\begin{document}

\section*{Proximal Gradient Method (Forward–Backward Splitting)}

\section*{Mathematical Formulation}
Let $f:\R^{d}\!\to\!\R$ be convex and $L$‑smooth (i.e. $\nabla f$ is $L$‑Lipschitz) and let $g:\R^{d}\!\to\!\R\cup\{+\infty\}$ be proper, closed and convex.  
Define the composite objective
\[
F(x)\;:=\;f(x)+g(x),\qquad x\in\R^{d}.
\]
For a stepsize $\eta>0$ the \emph{proximal gradient update} is
\begin{equation}
\boxed{%
x_{k+1}\;=\;\prox_{\eta g}\!\bigl(x_{k}-\eta\nabla f(x_{k})\bigr)},
\qquad k=0,1,2,\dots
\tag{1}
\end{equation}
where
\[
\prox_{\eta g}(u)\;:=\;\arg\min_{x\in\R^{d}}\Bigl\{ g(x)+\frac{1}{2\eta}\|x-u\|^{2}\Bigr\}.
\]

\section*{Pseudocode}
\begin{algorithm}[H]
\caption{Proximal Gradient Method}
\begin{algorithmic}[1]
\REQUIRE Initial point $x_{0}\in\R^{d}$, stepsize $\eta\in(0,2/L)$
\FOR{$k=0,1,2,\dots$}
    \STATE Compute gradient $g_{k}\gets\nabla f(x_{k})$
    \STATE Take a forward step $y_{k}\gets x_{k}-\eta g_{k}$
    \STATE Backward (prox) step $x_{k+1}\gets\prox_{\eta g}(y_{k})$
    \IF{termination criterion satisfied}
        \STATE \textbf{break}
    \ENDIF
\ENDFOR
\RETURN $x_{k+1}$
\end{algorithmic}
\end{algorithm}

\section*{Dry Run Example}
Consider the scalar problem
\[
f(w)=(w-3)^{2},\qquad g\equiv0,\qquad \eta=0.1,\qquad w_{0}=0.
\]
Since $g=0$, $\prox_{\eta g}(u)=u$.

\begin{center}
\begin{tabular}{c|c|c|c}
Iteration $k$ & $w_{k}$ & $\nabla f(w_{k})$ & $w_{k+1}=w_{k}-\eta\nabla f(w_{k})$ \\ \hline
0 & $0$ & $2(0-3)=-6$ & $0-0.1(-6)=0.6$ \\
1 & $0.6$ & $2(0.6-3)=-4.8$ & $0.6-0.1(-4.8)=1.08$ \\
2 & $1.08$ & $2(1.08-3)=-3.84$ & $1.08-0.1(-3.84)=1.464$ \\
\end{tabular}
\end{center}
Objective values:
\[
f(0)=9,\; f(0.6)=5.76,\; f(1.08)=3.6864,\; f(1.464)=2.3405.
\]
The iterates move monotonically toward the minimiser $w^{\star}=3$.

\newpage
\section*{Assumptions}
\begin{assumption}[Convexity and Smoothness of $f$]\label{as:convexsmooth}
$f$ is convex and differentiable with $L$‑Lipschitz gradient:
\[
\|\nabla f(x)-\nabla f(y)\|\le L\|x-y\|,\qquad\forall x,y\in\R^{d}.
\]
\end{assumption}

\begin{assumption}[Convexity of $g$]\label{as:gconvex}
$g$ is proper, closed and convex (possibly nonsmooth). Its proximal operator $\prox_{\eta g}$ is well defined for any $\eta>0$.
\end{assumption}

\begin{assumption}[Stepsize]\label{as:stepsize}
The stepsize satisfies
\[
0<\eta<\frac{2}{L}.
\]
\end{assumption}

\section*{Main Convergence Theorem}
\begin{theorem}[Ergodic $O(1/k)$ rate]\label{thm:rate}
Let $\{x_{k}\}$ be generated by \eqref{1} under Assumptions \ref{as:convexsmooth}–\ref{as:stepsize} and let $x^{\star}\in\arg\min F$. Then for all $K\ge1$,
\[
\frac{1}{K}\sum_{k=0}^{K-1}\bigl(F(x_{k})-F(x^{\star})\bigr)
\;\le\;
\frac{\|x_{0}-x^{\star}\|^{2}}{2\eta K}.
\tag{2}
\]
Consequently $F(x_{k})\to F(x^{\star})$ and the worst‑case objective gap decays as $O(1/k)$.
\end{theorem}

\section*{Theoretical Properties}
\begin{itemize}
\item \textbf{Stability.} Because $\prox_{\eta g}$ is non‑expansive (Lemma~\ref{lem:nonexp}), the iteration mapping $T_{\eta}(x)=\prox_{\eta g}(x-\eta\nabla f(x))$ is firmly non‑expansive for $0<\eta<2/L$, guaranteeing boundedness of the trajectory.
\item \textbf{Stepsize Sensitivity.} The bound \eqref{2} is monotone decreasing in $\eta$ up to the limit $2/L$; choosing $\eta$ close to $2/L$ yields the smallest constant $\frac{1}{2\eta}$.
\item \textbf{Comparison.} When $g\equiv0$, \eqref{1} reduces to gradient descent and \eqref{2} recovers the classic $O(1/k)$ rate. For nonsmooth $g$, the same rate holds without any extra assumptions beyond convexity.
\end{itemize}

\section*{Discussion}
The theorem shows that the proximal gradient method attains the optimal sublinear rate for the broad class of convex composite problems. The proof relies only on first‑order smoothness of $f$ and the non‑expansiveness of the proximal map; no stochasticity or strong convexity is required. Extensions to accelerated schemes or stochastic variants follow by strengthening Lemma~\ref{lem:nonexp} with momentum or variance‑reduction arguments.

\newpage
\section*{Preliminaries (Definitions and Lemmas)}
\begin{lemma}[Descent Lemma]\label{lem:descent}
If $\nabla f$ is $L$‑Lipschitz, then for any $x,y\in\R^{d}$,
\[
f(y)\le f(x)+\langle\nabla f(x),y-x\rangle+\frac{L}{2}\|y-x\|^{2}.
\tag{3}
\]
\end{lemma}
\begin{proof}
Standard consequence of the integral form of the gradient and the Lipschitz condition. \hfill $\square$
\end{proof}

\begin{lemma}[Non‑expansiveness of the proximal operator]\label{lem:nonexp}
For any $\eta>0$ and any $u,v\in\R^{d}$,
\[
\bigl\|\prox_{\eta g}(u)-\prox_{\eta g}(v)\bigr\|\le\|u-v\|.
\tag{4}
\]
\end{lemma}
\begin{proof}
Let $p=\prox_{\eta g}(u)$ and $q=\prox_{\eta g}(v)$. By optimality,
\[
0\in\partial g(p)+\frac{1}{\eta}(p-u),\qquad
0\in\partial g(q)+\frac{1}{\eta}(q-v).
\]
Subtracting the two inclusions and using monotonicity of $\partial g$ yields
\[
\langle p-q-\!(u-v),\,p-q\rangle\ge0
\;\Longrightarrow\;
\|p-q\|^{2}\le\langle u-v,\,p-q\rangle\le\|u-v\|\|p-q\|,
\]
which gives (4). \hfill $\square$
\end{proof}

\section*{Main Proof (Step‑by‑Step Derivation)}
We prove Theorem~\ref{thm:rate}.  
Let $x^{\star}\in\arg\min F$ and define the forward‑backward operator
\[
T_{\eta}(x):=\prox_{\eta g}\bigl(x-\eta\nabla f(x)\bigr).
\]
The iteration is $x_{k+1}=T_{\eta}(x_{k})$.

\noindent\textbf{[Step 1]} (Optimality condition of the proximal step).  
By definition of $\prox_{\eta g}$,
\[
x_{k+1}= \prox_{\eta g}\bigl(x_{k}-\eta\nabla f(x_{k})\bigr)
\Longleftrightarrow
0\in\partial g(x_{k+1})+\frac{1}{\eta}\bigl(x_{k+1}-x_{k}+\eta\nabla f(x_{k})\bigr).
\tag{5}
\]

\noindent\textbf{[Step 2]} (Variational inequality).  
For any $z\in\R^{d}$, multiply (5) by $z-x_{k+1}$ and use monotonicity of $\partial g$:
\[
\big\langle \tfrac{1}{\eta}(x_{k}-x_{k+1})+\nabla f(x_{k}),\,z-x_{k+1}\big\rangle
\ge g(z)-g(x_{k+1}).
\tag{6}
\]

\noindent\textbf{[Step 3]} (Choose $z=x^{\star}$).  
Since $x^{\star}$ minimizes $F$, we have $0\in\partial g(x^{\star})+\nabla f(x^{\star})$. Adding the inequality obtained by substituting $z=x^{\star}$ in (6) and the analogous inequality with $x^{\star}$ gives
\[
\big\langle \tfrac{1}{\eta}(x_{k}-x_{k+1})+\nabla f(x_{k})-\nabla f(x^{\star}),\,x^{\star}-x_{k+1}\big\rangle
\ge F(x^{\star})-F(x_{k+1}).
\tag{7}
\]

\noindent\textbf{[Step 4]} (Expand the inner product).  
Rewrite the left‑hand side of (7) as
\[
\frac{1}{\eta}\langle x_{k}-x_{k+1},x^{\star}-x_{k+1}\rangle
+\langle\nabla f(x_{k})-\nabla f(x^{\star}),x^{\star}-x_{k+1}\rangle .
\tag{8}
\]

\noindent\textbf{[Step 5]} (Algebraic identity).  
For any vectors $a,b$,
\[
\langle a-b,a-c\rangle=\tfrac12\bigl(\|a-c\|^{2}+\|b-c\|^{2}-\|a-b\|^{2}\bigr).
\tag{9}
\]
Apply (9) with $a=x_{k}$, $b=x_{k+1}$, $c=x^{\star}$ to the first term of (8):
\[
\frac{1}{\eta}\langle x_{k}-x_{k+1},x^{\star}-x_{k+1}\rangle
=\frac{1}{2\eta}\Bigl(\|x_{k}-x^{\star}\|^{2}
+\|x_{k+1}-x^{\star}\|^{2}
-\|x_{k}-x_{k+1}\|^{2}\Bigr).
\tag{10}
\]

\noindent\textbf{[Step 6]} (Use convexity of $f$).  
Convexity of $f$ yields
\[
f(x^{\star})\ge f(x_{k})+\langle\nabla f(x_{k}),x^{\star}-x_{k}\rangle,
\qquad
f(x_{k+1})\ge f(x_{k})+\langle\nabla f(x_{k}),x_{k+1}-x_{k}\rangle.
\tag{11}
\]
Subtract the second inequality from the first and rearrange:
\[
\langle\nabla f(x_{k})-\nabla f(x^{\star}),x^{\star}-x_{k+1}\rangle
\le f(x^{\star})-f(x_{k+1}) .
\tag{12}
\]

\noindent\textbf{[Step 7]} (Combine (10) and (12) in (8)).  
Plugging (10) and (12) into (8) we obtain
\[
\frac{1}{2\eta}\Bigl(\|x_{k}-x^{\star}\|^{2}
+\|x_{k+1}-x^{\star}\|^{2}
-\|x_{k}-x_{k+1}\|^{2}\Bigr)
+f(x^{\star})-f(x_{k+1})
\ge F(x^{\star})-F(x_{k+1}).
\tag{13}
\]

\noindent\textbf{[Step 8]} (Cancel $f(x^{\star})$).  
Since $F(x)=f(x)+g(x)$, the term $f(x^{\star})$ appears on both sides of (13) and cancels, leaving
\[
\frac{1}{2\eta}\Bigl(\|x_{k}-x^{\star}\|^{2}
-\|x_{k+1}-x^{\star}\|^{2}
-\|x_{k}-x_{k+1}\|^{2}\Bigr)
\ge g(x_{k+1})-g(x^{\star}).
\tag{14}
\]

\noindent\textbf{[Step 9]} (Form the fundamental inequality).  
Rearranging (14) gives
\[
F(x_{k+1})-F(x^{\star})
\le \frac{1}{2\eta}\Bigl(\|x_{k}-x^{\star}\|^{2}
-\|x_{k+1}-x^{\star}\|^{2}\Bigr)
-\frac{1}{2\eta}\|x_{k}-x_{k+1}\|^{2}.
\tag{15}
\]

\noindent\textbf{[Step 10]} (Sum over $k=0,\dots,K-1$).  
Summing (15) telescopically yields
\[
\sum_{k=0}^{K-1}\bigl(F(x_{k+1})-F(x^{\star})\bigr)
\le \frac{1}{2\eta}\|x_{0}-x^{\star}\|^{2}
-\frac{1}{2\eta}\sum_{k=0}^{K-1}\|x_{k}-x_{k+1}\|^{2}.
\tag{16}
\]

\noindent\textbf{[Step 11]} (Discard the non‑negative term).  
Since $\|x_{k}-x_{k+1}\|^{2}\ge0$, we drop the last sum to obtain the simpler bound
\[
\sum_{k=0}^{K-1}\bigl(F(x_{k+1})-F(x^{\star})\bigr)
\le \frac{1}{2\eta}\|x_{0}-x^{\star}\|^{2}.
\tag{17}
\]

\noindent\textbf{[Step 12]} (Averaging).  
Divide both sides of (17) by $K$ and relabel the index ($k\gets k-1$) to reach
\[
\frac{1}{K}\sum_{k=0}^{K-1}\bigl(F(x_{k})-F(x^{\star})\bigr)
\le \frac{\|x_{0}-x^{\star}\|^{2}}{2\eta K},
\]
which is exactly the claim \eqref{2}.  $\square$

\section*{Rate Derivation (With Constants)}
From \eqref{2} we have the explicit constant
\[
C(\eta):=\frac{\|x_{0}-x^{\star}\|^{2}}{2\eta},
\qquad\text{so that}\qquad
\frac{1}{K}\sum_{k=0}^{K-1}\bigl(F(x_{k})-F^{\star}\bigr)\le\frac{C(\eta)}{K}.
\]
The bound is monotone decreasing in $\eta$ on $(0,2/L)$; the optimal admissible stepsize is $\eta^{\star}=2/L$, giving
\[
C(\eta^{\star})=\frac{L}{4}\|x_{0}-x^{\star}\|^{2}.
\]

\section*{Special Cases}
\begin{itemize}
\item \textbf{Pure gradient descent} ($g\equiv0$).  The proximal map reduces to the identity, and inequality \eqref{15} becomes the classic gradient‑descent bound.
\item \textbf{Indicator of a convex set $\mathcal{C}$}.  If $g=\iota_{\mathcal{C}}$, then $\prox_{\eta g}$ is the Euclidean projection onto $\mathcal{C}$.  The same $O(1/k)$ rate holds for projected gradient descent.
\item \textbf{Strongly convex $F$}.  If $F$ is $\mu$‑strongly convex, a standard argument upgrades (15) to a linear rate $O\bigl((1-\eta\mu)^{k}\bigr)$.
\end{itemize}

\end{document}